% Options for packages loaded elsewhere
\PassOptionsToPackage{unicode}{hyperref}
\PassOptionsToPackage{hyphens}{url}
\documentclass[
]{article}
\usepackage{xcolor}
\usepackage[margin=1in]{geometry}
\usepackage{amsmath,amssymb}
\setcounter{secnumdepth}{-\maxdimen} % remove section numbering
\usepackage{iftex}
\ifPDFTeX
  \usepackage[T1]{fontenc}
  \usepackage[utf8]{inputenc}
  \usepackage{textcomp} % provide euro and other symbols
\else % if luatex or xetex
  \usepackage{unicode-math} % this also loads fontspec
  \defaultfontfeatures{Scale=MatchLowercase}
  \defaultfontfeatures[\rmfamily]{Ligatures=TeX,Scale=1}
\fi
\usepackage{lmodern}
\ifPDFTeX\else
  % xetex/luatex font selection
\fi
% Use upquote if available, for straight quotes in verbatim environments
\IfFileExists{upquote.sty}{\usepackage{upquote}}{}
\IfFileExists{microtype.sty}{% use microtype if available
  \usepackage[]{microtype}
  \UseMicrotypeSet[protrusion]{basicmath} % disable protrusion for tt fonts
}{}
\makeatletter
\@ifundefined{KOMAClassName}{% if non-KOMA class
  \IfFileExists{parskip.sty}{%
    \usepackage{parskip}
  }{% else
    \setlength{\parindent}{0pt}
    \setlength{\parskip}{6pt plus 2pt minus 1pt}}
}{% if KOMA class
  \KOMAoptions{parskip=half}}
\makeatother
\usepackage{color}
\usepackage{fancyvrb}
\newcommand{\VerbBar}{|}
\newcommand{\VERB}{\Verb[commandchars=\\\{\}]}
\DefineVerbatimEnvironment{Highlighting}{Verbatim}{commandchars=\\\{\}}
% Add ',fontsize=\small' for more characters per line
\usepackage{framed}
\definecolor{shadecolor}{RGB}{248,248,248}
\newenvironment{Shaded}{\begin{snugshade}}{\end{snugshade}}
\newcommand{\AlertTok}[1]{\textcolor[rgb]{0.94,0.16,0.16}{#1}}
\newcommand{\AnnotationTok}[1]{\textcolor[rgb]{0.56,0.35,0.01}{\textbf{\textit{#1}}}}
\newcommand{\AttributeTok}[1]{\textcolor[rgb]{0.13,0.29,0.53}{#1}}
\newcommand{\BaseNTok}[1]{\textcolor[rgb]{0.00,0.00,0.81}{#1}}
\newcommand{\BuiltInTok}[1]{#1}
\newcommand{\CharTok}[1]{\textcolor[rgb]{0.31,0.60,0.02}{#1}}
\newcommand{\CommentTok}[1]{\textcolor[rgb]{0.56,0.35,0.01}{\textit{#1}}}
\newcommand{\CommentVarTok}[1]{\textcolor[rgb]{0.56,0.35,0.01}{\textbf{\textit{#1}}}}
\newcommand{\ConstantTok}[1]{\textcolor[rgb]{0.56,0.35,0.01}{#1}}
\newcommand{\ControlFlowTok}[1]{\textcolor[rgb]{0.13,0.29,0.53}{\textbf{#1}}}
\newcommand{\DataTypeTok}[1]{\textcolor[rgb]{0.13,0.29,0.53}{#1}}
\newcommand{\DecValTok}[1]{\textcolor[rgb]{0.00,0.00,0.81}{#1}}
\newcommand{\DocumentationTok}[1]{\textcolor[rgb]{0.56,0.35,0.01}{\textbf{\textit{#1}}}}
\newcommand{\ErrorTok}[1]{\textcolor[rgb]{0.64,0.00,0.00}{\textbf{#1}}}
\newcommand{\ExtensionTok}[1]{#1}
\newcommand{\FloatTok}[1]{\textcolor[rgb]{0.00,0.00,0.81}{#1}}
\newcommand{\FunctionTok}[1]{\textcolor[rgb]{0.13,0.29,0.53}{\textbf{#1}}}
\newcommand{\ImportTok}[1]{#1}
\newcommand{\InformationTok}[1]{\textcolor[rgb]{0.56,0.35,0.01}{\textbf{\textit{#1}}}}
\newcommand{\KeywordTok}[1]{\textcolor[rgb]{0.13,0.29,0.53}{\textbf{#1}}}
\newcommand{\NormalTok}[1]{#1}
\newcommand{\OperatorTok}[1]{\textcolor[rgb]{0.81,0.36,0.00}{\textbf{#1}}}
\newcommand{\OtherTok}[1]{\textcolor[rgb]{0.56,0.35,0.01}{#1}}
\newcommand{\PreprocessorTok}[1]{\textcolor[rgb]{0.56,0.35,0.01}{\textit{#1}}}
\newcommand{\RegionMarkerTok}[1]{#1}
\newcommand{\SpecialCharTok}[1]{\textcolor[rgb]{0.81,0.36,0.00}{\textbf{#1}}}
\newcommand{\SpecialStringTok}[1]{\textcolor[rgb]{0.31,0.60,0.02}{#1}}
\newcommand{\StringTok}[1]{\textcolor[rgb]{0.31,0.60,0.02}{#1}}
\newcommand{\VariableTok}[1]{\textcolor[rgb]{0.00,0.00,0.00}{#1}}
\newcommand{\VerbatimStringTok}[1]{\textcolor[rgb]{0.31,0.60,0.02}{#1}}
\newcommand{\WarningTok}[1]{\textcolor[rgb]{0.56,0.35,0.01}{\textbf{\textit{#1}}}}
\usepackage{longtable,booktabs,array}
\usepackage{calc} % for calculating minipage widths
% Correct order of tables after \paragraph or \subparagraph
\usepackage{etoolbox}
\makeatletter
\patchcmd\longtable{\par}{\if@noskipsec\mbox{}\fi\par}{}{}
\makeatother
% Allow footnotes in longtable head/foot
\IfFileExists{footnotehyper.sty}{\usepackage{footnotehyper}}{\usepackage{footnote}}
\makesavenoteenv{longtable}
\usepackage{graphicx}
\makeatletter
\newsavebox\pandoc@box
\newcommand*\pandocbounded[1]{% scales image to fit in text height/width
  \sbox\pandoc@box{#1}%
  \Gscale@div\@tempa{\textheight}{\dimexpr\ht\pandoc@box+\dp\pandoc@box\relax}%
  \Gscale@div\@tempb{\linewidth}{\wd\pandoc@box}%
  \ifdim\@tempb\p@<\@tempa\p@\let\@tempa\@tempb\fi% select the smaller of both
  \ifdim\@tempa\p@<\p@\scalebox{\@tempa}{\usebox\pandoc@box}%
  \else\usebox{\pandoc@box}%
  \fi%
}
% Set default figure placement to htbp
\def\fps@figure{htbp}
\makeatother
\setlength{\emergencystretch}{3em} % prevent overfull lines
\providecommand{\tightlist}{%
  \setlength{\itemsep}{0pt}\setlength{\parskip}{0pt}}
\usepackage[]{natbib}
\bibliographystyle{plainnat}
\usepackage{bookmark}
\IfFileExists{xurl.sty}{\usepackage{xurl}}{} % add URL line breaks if available
\urlstyle{same}
\hypersetup{
  pdftitle={aucmat: matrix-first ROC analysis for high-dimensional biomarker screening in R},
  hidelinks,
  pdfcreator={LaTeX via pandoc}}

\title{aucmat: matrix-first ROC analysis for high-dimensional biomarker
screening in R}
\author{true}
\date{}

\begin{document}
\maketitle
\begin{abstract}
\textbf{Summary:} The area under the receiver operating characteristic
curve (AUC) is the standard metric for evaluating continuous biomarkers
against binary outcomes. When screening thousands of molecular
biomarkers, researchers need matrix-wide AUC estimation with proper
inference, multiplicity adjustment, and stability assessment. We present
aucmat, an R package providing a matrix-first architecture:
\texttt{aucmat(X,\ y)} computes direction-preserving AUCs with DeLong or
stratified-bootstrap confidence intervals for every column of a numeric
matrix in a single call. The package supports Benjamini-Hochberg
multiplicity adjustment, three hypothesis-testing frameworks for paired
comparisons (superiority, non-inferiority, equivalence via TOST), an
omnibus Wald test for equality of multiple correlated AUCs via the joint
DeLong covariance matrix, and bootstrap rank-stability analysis.
Cross-validated screening eliminates same-data optimism. A two-stage
Hurdle-AUC model handles zero-inflated data (scRNA-seq, microbiome)
where standard AUC degrades to near chance. Eight simulation generators
and a validation harness produce correlated biomarkers with controlled
properties under parametrized correlation structures with built-in
calibration diagnostics. The package provides over a dozen visualization
functions including ranked discrimination plots, volcano plots, forest
plots, ROC curves with confidence ribbons, precision-recall curves, and
Hurdle-AUC diagnostics. On a typical proteomics screen, aucmat avoids
the per-column looping required by single-biomarker tools and provides
complete inference in a single function call. The package is freely
available under the MIT license at
\url{https://github.com/vanhungtran/aucmat}. \textbf{Availability and
Implementation:} R package (R \textgreater= 4.0.0), MIT license. Source:
\url{https://github.com/vanhungtran/aucmat}. Documentation:
\url{https://vanhungtran.github.io/aucmat}. Dependencies: pROC
(\textgreater= 1.18), mvtnorm, ggplot2, rlang. \textbf{Contact:}
\href{mailto:lucas.tran@ck-care.ch}{\nolinkurl{lucas.tran@ck-care.ch}}
\end{abstract}

\section{Introduction}\label{introduction}

Receiver operating characteristic (ROC) analysis is the dominant
framework for evaluating continuous biomarkers against binary outcomes
in diagnostic medicine, prognostic modeling, and molecular screening
\citep{Hanley1982, DeLong1988, Pepe2004}. The area under the ROC curve
(AUC) summarizes the probability that a randomly chosen positive subject
has a higher biomarker value than a randomly chosen negative subject.
Its interpretability -- 0.5 for chance, 1.0 for perfect separation --
has made it the default metric in omics-scale biomarker discovery
\citep{Robin2011}.

Modern molecular profiling technologies routinely measure thousands of
features on a few hundred subjects. A proteomics panel may contain 5,000
proteins; a metabolomics screen 2,000 metabolites; a transcriptomics
experiment 20,000 genes. The core statistical task -- ``which biomarkers
discriminate cases from controls?'' -- requires computing and testing an
AUC for each feature, adjusting for multiplicity, comparing top
candidates, and assessing the stability of the rankings under sampling
variability. Despite decades of methodological development, no
widely-used R package integrates these steps into a single, reproducible
workflow.

The existing R ecosystem provides excellent point solutions.
\textbf{pROC} \citep{Robin2011} is the gold standard for
single-biomarker ROC inference: it computes DeLong confidence intervals
using the Sun \& Xu (2014) O(N log N) algorithm \citep{Sun2014}, tests
paired or unpaired AUC differences via \texttt{roc.test()}, handles
partial AUC with bootstrap inference, provides binormal and
density-smoothed ROC curves, computes covariance between paired AUCs via
\texttt{cov()}, and supports multiclass AUC via Hand \& Till
\citep{Hand2001}. However, pROC operates one biomarker at a time --
screening a matrix of 10,000 features requires an explicit per-column
loop. \textbf{ROCR} \citep{Sing2005} offers flexible visualization with
28 performance measures and cross-validation averaging but provides no
statistical inference (no confidence intervals, no hypothesis tests for
AUC differences). \textbf{precrec} \citep{Saito2017} adds
precision-recall curves and fast C++-based AUC computation, with
confidence bands for multiple test sets and functions for partial AUC
(\texttt{part()}, \texttt{pauc()}) and AUC confidence intervals
(\texttt{auc\_ci()}), but provides no DeLong inference, no multiplicity
adjustment, and no paired comparison tests. \textbf{colAUC} from caTools
\citep{Tuszynski2024} computes column-wise AUCs efficiently via the
Wilcoxon rank-sum test and supports multiclass AUC via column means --
but provides no inference whatsoever (no standard errors, confidence
intervals, or p-values). \textbf{dtComb} \citep{Yerlitas2025} provides
142 methods for combining two diagnostic markers into a composite score,
with 30+ cutoff optimization methods, resampling, and caret integration
-- but is limited to exactly two biomarkers and does not perform
matrix-wide screening. \textbf{cancerclass} \citep{Budczies2024} builds
validated classifiers from high-dimensional molecular data using
nearest-centroid prediction with multiple random validation, but does
not offer DeLong-based screening or bootstrap stability analysis.

aucmat fills this gap. Its design philosophy is \emph{matrix-first}: the
primary function \texttt{aucmat(X,\ y)} accepts a numeric matrix
\(\mathbf{X} \in
\mathbb{R}^{n \times p}\) and a binary outcome vector \(\mathbf{y} \in
\{0,1\}^n\), and returns a complete screening table with point
estimates, standard errors, confidence intervals, raw and
multiplicity-adjusted p-values, and diagnostic metadata for every
biomarker in a single call. Downstream functions for comparison,
stability, cross-validation, panel construction, and power analysis
consume the same S3 object, preserving outcome encoding and inferential
settings throughout.

\section{Methods}\label{methods}

\subsection{Direction-preserving AUC}\label{direction-preserving-auc}

For biomarker \(X\) and binary outcome \(Y \in \{0,1\}\), the AUC equals
the Mann-Whitney U probability:
\(\text{AUC} = P(X_i > X_j \mid Y_i = 1, Y_j =
0)\) \citep{Hanley1982}. Most ROC software reports
\(\max(\text{AUC}, 1 -
\text{AUC})\), silently reversing the biomarker to force the reported
value above 0.5. aucmat deliberately preserves the observed direction,
reporting three values per biomarker: \(\text{AUC}_{\text{raw}}\) (the
Mann-Whitney estimate on the original numeric scale, which may be below
0.5),
\(\text{AUC}_{\text{strength}} = 0.5 + |\text{AUC}_{\text{raw}} - 0.5|\)
(for ranking), and \texttt{effect\_direction}
(\texttt{higher\_in\_positive} or \texttt{lower\_in\_positive}). This
distinction matters biologically: a tumor suppressor gene with
\(\text{AUC}_{\text{raw}} = 0.15\) discriminates just as strongly as an
oncogene with \(\text{AUC}_{\text{raw}} = 0.85\), but with the opposite
biological interpretation.

\subsection{DeLong variance and joint
covariance}\label{delong-variance-and-joint-covariance}

The DeLong variance estimator \citep{DeLong1988, Sun2014} uses placement
values
\(V_{10}(X_i) = n_0^{-1} \sum_{j:Y_j=0} [\mathbf{1}(X_j < X_i) + 0.5 \cdot
\mathbf{1}(X_j = X_i)]\) for positive subjects and
\(V_{01}(X_j) = n_1^{-1}
\sum_{i:Y_i=1} [\mathbf{1}(X_i > X_j) + 0.5 \cdot \mathbf{1}(X_i = X_j)]\)
for negative subjects. Let
\(S_{10} = (n_1 - 1)^{-1} \sum_{i \in \mathcal{I}_1}
(V_{10}(X_i) - \bar{V}_{10})^2\) and
\(S_{01} = (n_0 - 1)^{-1} \sum_{j \in
\mathcal{I}_0} (V_{01}(X_j) - \bar{V}_{01})^2\) be the sample variances
of the placement values. The DeLong variance estimator is
\(\widehat{\text{Var}}(\widehat{\text{AUC}}) = S_{10}/n_1 + S_{01}/n_0\).
For \(p\) biomarkers on common subjects, the \(p \times p\) joint
covariance matrix
\(\widehat{\mathbf{\Sigma}} = n_1^{-1}\widehat{\text{Cov}}(\mathbf{V}_{10})
+ n_0^{-1}\widehat{\text{Cov}}(\mathbf{V}_{01})\) supports an omnibus
Wald test of \(H_0: \text{AUC}_1 = \cdots = \text{AUC}_p\). Let
\(\mathbf{C}\) be a \((p-1) \times p\) contrast matrix (default:
comparing biomarkers \(2,\ldots,p\) against biomarker 1). The Wald
statistic \(W = (\mathbf{C}\hat{\boldsymbol{\theta}})^\top
(\mathbf{C}\widehat{\mathbf{\Sigma}}\mathbf{C}^\top)^{-1}
(\mathbf{C}\hat{\boldsymbol{\theta}}) \sim \chi^2_{r}\) uses the
effective rank \(r\) of the contrast covariance matrix, determined via
eigenvalue tolerance. A generalized inverse is used on the stable
subspace only; rank-zero or numerically unstable cases return
\texttt{status\ =\ "non\_testable\_singular"}. For \(p = 2\), \(W\)
equals the squared paired-DeLong \(z\)-statistic within numerical
tolerance, serving as an internal consistency check.

\subsection{Multiplicity adjustment and
stability}\label{multiplicity-adjustment-and-stability}

Screening \(p\) biomarkers produces \(p\) hypothesis tests. aucmat
applies Benjamini-Hochberg \citep{Benjamini1995}, Holm, or Bonferroni
correction via \texttt{stats::p.adjust()}, reporting adjusted p-values
as \(q\)-values. For equivalence (TOST) \citep{Schuirmann1987}, the
reported \(p_{\text{equiv}} =
\max(p_{\text{lower}}, p_{\text{upper}})\) and multiplicity adjustment
applies to this combined value.

Bootstrap rank-stability analysis (\texttt{auc\_stability()}) resamples
the data \(B\) times with replacement (stratified by outcome),
recomputes the full screening for each replicate, and reports median/IQR
rank distributions, top-\(k\) selection probabilities, and pairwise
co-selection frequencies -- quantifying how reproducible the ranking is
under sampling variability \citep{Efron1994}.

\subsection{Hurdle-AUC for zero-inflated
data}\label{hurdle-auc-for-zero-inflated-data}

Standard AUC assumes continuous biomarkers following a shifted normal
distribution per class. This assumption breaks on zero-inflated data
(scRNA-seq, microbiome, detection-limited assays) where 60\%+ of values
are exactly zero. Under the standard model, zeros would need to arise
from the tail of a normal distribution below zero -- biologically
implausible. The Hurdle-AUC uses a two-stage model:

\textbf{Stage 1 -- Zero hurdle:}
\(P(X_j = 0 \mid Y) = \text{logistic}(\beta_0 +
\beta_1 Y)\) models the binary expression decision via logistic
regression.

\textbf{Stage 2 -- Magnitude:} Standard AUC on the non-zero values
captures magnitude discrimination among expressed observations.

The composite score \(S_i\) for observation \(i\) is:
\(S_i = \hat{p}_i\) if \(X_{ij} = 0\), and
\(S_i = 0.5 \cdot \hat{p}_i + 0.5 \cdot
\tilde{x}_{ij}\) if \(X_{ij} > 0\), where \(\hat{p}_i\) is the predicted
probability of being non-zero from Stage 1, and \(\tilde{x}_{ij}\) is
the normalized biomarker value.

Simulation confirms recovery from standard AUC \(\approx\) 0.53 to
Hurdle-AUC \(\approx\) 0.90 on zero-inflated biomarkers by acknowledging
that ``expression on/off'' and ``expression level when on'' are separate
biological processes.

\subsection{Multivariate simulation}\label{multivariate-simulation}

Under the equal-variance binormal model,
\(X_k \mid Y \sim \mathcal{N}(\mu_{kY},
\sigma^2)\) with
\(\delta_k = (\mu_{1k} - \mu_{0k})/\sigma = \sqrt{2} \cdot
\Phi^{-1}(\text{AUC}_k)\). For \(p\) biomarkers with correlation matrix
\(\mathbf{R}\), the class-conditional distribution is
\(\mathbf{X} \mid Y \sim
\mathcal{N}_p(\boldsymbol{\mu}_Y, \mathbf{R})\), where
\(\boldsymbol{\mu}_0\) and \(\boldsymbol{\mu}_1\) are centred so the
\(n\)-weighted grand mean is zero. \texttt{simulate\_auc\_matrix()}
generates data under this model with four parametrized correlation
structures (user, exchangeable, AR(1), block) and feasibility
diagnostics with optional nearest positive-definite projection.

\texttt{simulate\_auc\_copula()} uses a two-phase approach: Phase 1
applies a Gaussian copula (rank-transform to uniform, probit to normal
scores) that preserves the target correlation structure; Phase 2 applies
iterative mean perturbation with decaying step size to fine-tune each
biomarker's AUC. This separates correlation control (copula) from
discriminative signal (perturbation), avoiding the AUC-correlation
tradeoff inherent in sequential binormal methods.

\texttt{simulate\_hurdle\_auc()} implements the two-stage hurdle model:
Bernoulli zeros with class-specific rates, and log-normal magnitudes for
expressed values with binormal-calibrated mean separation, clipped to
\([0, \infty)\).

\texttt{validate\_simulation()} repeats a specification across
independent draws and reports bias, RMSE, Monte Carlo standard error,
and target-interval hit rates for both AUCs and pairwise correlations --
a single draw is never evidence that a simulator is calibrated.

\section{Implementation}\label{implementation}

aucmat is implemented in pure R (\textasciitilde5,700 lines across 35
source files) with minimized dependencies. The primary function
\texttt{aucmat()} performs six steps in sequence: (1) input validation
(matrix dimensions, column names, outcome levels, infinite values); (2)
outcome normalization via \texttt{.normalize\_binary\_y()} (accepting
logical, numeric 0/1, factor, or character outcomes); (3) matrix-wide
AUC computation via the Mann-Whitney rank formulation; (4) DeLong or
bootstrap inference via \texttt{.apply\_inference()} with configurable
alternative hypothesis (\texttt{two.sided}, \texttt{greater},
\texttt{less}); (5) multiplicity adjustment via
\texttt{.adjust\_pvalues()}; and (6) result assembly with ranking by
discrimination strength and feature-level warnings for small class
counts or high missingness.

All primary functions return S3 objects with \texttt{print()},
\texttt{summary()}, \texttt{plot()}, and (where appropriate)
\texttt{as.data.frame()} and \texttt{subset()} methods. The package
stores outcome encoding, positive-class designation, and sample masks in
the returned object, eliminating the need for downstream functions to
re-infer factor direction or re-derive class counts.

Structured error classes (inheriting from \texttt{error}) allow
programmatic handling of common failures:
\texttt{aucmat\_invalid\_outcome},
\texttt{aucmat\_invalid\_correlation},
\texttt{aucmat\_infeasible\_targets},
\texttt{aucmat\_insufficient\_sample}, and
\texttt{aucmat\_resampling\_failure}. An RNG-safety mechanism
(\texttt{with\_seed()}) saves the global \texttt{.Random.seed}, executes
the computation, and restores the previous state on exit, including when
\texttt{.Random.seed} was previously absent and when an error occurs.

The package imports pROC (ROC curve objects and smoothed ROC), mvtnorm
(multivariate normal generation), ggplot2 (visualization), and rlang
(tidy evaluation). Optional Suggests include GGally (pairs plots for
simulated data), glmnet (penalized panel models), matrixStats (fast
column-wise ranks), knitr and rmarkdown (vignette building), and ggrepel
(non-overlapping plot labels). The complete source, documentation, and
four vignettes are available at
\url{https://vanhungtran.github.io/aucmat}.

\section{Usage}\label{usage}

\subsection{Basic screening}\label{basic-screening}

\begin{Shaded}
\begin{Highlighting}[]
\FunctionTok{library}\NormalTok{(aucmat)}

\NormalTok{sim }\OtherTok{\textless{}{-}} \FunctionTok{simulate\_auc\_matrix}\NormalTok{(}
  \AttributeTok{n =} \DecValTok{500}\NormalTok{, }\AttributeTok{prevalence =} \FloatTok{0.3}\NormalTok{,}
  \AttributeTok{target\_aucs =} \FunctionTok{c}\NormalTok{(}\FloatTok{0.9}\NormalTok{, }\FloatTok{0.8}\NormalTok{, }\FloatTok{0.7}\NormalTok{, }\FloatTok{0.65}\NormalTok{),}
  \AttributeTok{correlation =} \FloatTok{0.3}\NormalTok{, }\AttributeTok{structure =} \StringTok{"exchangeable"}
\NormalTok{)}
\NormalTok{X }\OtherTok{\textless{}{-}} \FunctionTok{as.matrix}\NormalTok{(sim}\SpecialCharTok{$}\NormalTok{data[, }\DecValTok{1}\SpecialCharTok{:}\DecValTok{4}\NormalTok{])}
\NormalTok{y }\OtherTok{\textless{}{-}}\NormalTok{ sim}\SpecialCharTok{$}\NormalTok{data}\SpecialCharTok{$}\NormalTok{truth}

\NormalTok{fit }\OtherTok{\textless{}{-}} \FunctionTok{aucmat}\NormalTok{(X, y, }\AttributeTok{ci =} \StringTok{"delong"}\NormalTok{, }\AttributeTok{adjust =} \StringTok{"BH"}\NormalTok{)}
\FunctionTok{print}\NormalTok{(fit)}
\FunctionTok{summary}\NormalTok{(fit)}
\FunctionTok{plot\_auc\_rank}\NormalTok{(fit, }\AttributeTok{n\_label =} \DecValTok{4}\NormalTok{)}
\end{Highlighting}
\end{Shaded}

\subsection{Paired comparisons with hypothesis
testing}\label{paired-comparisons-with-hypothesis-testing}

\begin{Shaded}
\begin{Highlighting}[]
\CommentTok{\# Two{-}sided superiority (default: H0: delta = 0)}
\FunctionTok{compare\_auc}\NormalTok{(fit, X, y, }\AttributeTok{reference =} \StringTok{"X1"}\NormalTok{)}

\CommentTok{\# Non{-}inferiority: H0: AUC\_a \textless{}= AUC\_b {-} margin}
\FunctionTok{compare\_auc}\NormalTok{(fit, X, y, }\AttributeTok{reference =} \StringTok{"X1"}\NormalTok{,}
  \AttributeTok{hypothesis =} \StringTok{"noninferiority"}\NormalTok{, }\AttributeTok{margin =} \FloatTok{0.05}\NormalTok{)}

\CommentTok{\# Equivalence (TOST): H0: |AUC\_a {-} AUC\_b| \textgreater{}= margin}
\FunctionTok{compare\_auc}\NormalTok{(fit, X, y, }\AttributeTok{reference =} \StringTok{"X1"}\NormalTok{,}
  \AttributeTok{hypothesis =} \StringTok{"equivalence"}\NormalTok{, }\AttributeTok{margin =} \FloatTok{0.15}\NormalTok{, }\AttributeTok{adjust =} \StringTok{"BH"}\NormalTok{)}

\CommentTok{\# Omnibus Wald test: H0: AUC\_1 = AUC\_2 = ... = AUC\_p}
\FunctionTok{compare\_auc\_global}\NormalTok{(fit, X, y, }\AttributeTok{biomarkers =} \FunctionTok{paste0}\NormalTok{(}\StringTok{"X"}\NormalTok{, }\DecValTok{1}\SpecialCharTok{:}\DecValTok{4}\NormalTok{))}

\CommentTok{\# Single{-}biomarker DeLong test against chance}
\FunctionTok{roc\_test}\NormalTok{(X[, }\DecValTok{1}\NormalTok{], y)}
\end{Highlighting}
\end{Shaded}

\subsection{Hurdle-AUC for zero-inflated
data}\label{hurdle-auc-for-zero-inflated-data-1}

\begin{Shaded}
\begin{Highlighting}[]
\CommentTok{\# scRNA{-}seq: 55\% zeros in controls, 25\% in cases}
\NormalTok{sim }\OtherTok{\textless{}{-}} \FunctionTok{simulate\_hurdle\_auc}\NormalTok{(}
  \AttributeTok{n =} \DecValTok{400}\NormalTok{, }\AttributeTok{prevalence =} \FloatTok{0.3}\NormalTok{,}
  \AttributeTok{target\_hurdle\_aucs =} \FunctionTok{c}\NormalTok{(}\FloatTok{0.85}\NormalTok{, }\FloatTok{0.72}\NormalTok{),}
  \AttributeTok{zero\_rate\_neg =} \FunctionTok{c}\NormalTok{(}\FloatTok{0.55}\NormalTok{, }\FloatTok{0.80}\NormalTok{),}
  \AttributeTok{zero\_rate\_pos =} \FunctionTok{c}\NormalTok{(}\FloatTok{0.25}\NormalTok{, }\FloatTok{0.70}\NormalTok{)}
\NormalTok{)}
\NormalTok{fit\_hur }\OtherTok{\textless{}{-}} \FunctionTok{hurdle\_auc}\NormalTok{(}\FunctionTok{as.matrix}\NormalTok{(sim}\SpecialCharTok{$}\NormalTok{data[, }\DecValTok{1}\SpecialCharTok{:}\DecValTok{2}\NormalTok{]), sim}\SpecialCharTok{$}\NormalTok{data}\SpecialCharTok{$}\NormalTok{truth)}
\FunctionTok{print}\NormalTok{(fit\_hur)}
\CommentTok{\# Standard AUC \textasciitilde{}0.53 recovers to Hurdle AUC \textasciitilde{}0.90}
\end{Highlighting}
\end{Shaded}

\subsection{Cross-validation, panel construction, power
analysis}\label{cross-validation-panel-construction-power-analysis}

\begin{Shaded}
\begin{Highlighting}[]
\CommentTok{\# Cross{-}validated AUC (honest out{-}of{-}fold estimates)}
\NormalTok{cv }\OtherTok{\textless{}{-}} \FunctionTok{cv\_aucmat}\NormalTok{(X, y, }\AttributeTok{n\_folds =} \DecValTok{5}\NormalTok{, }\AttributeTok{seed =} \DecValTok{42}\NormalTok{)}

\CommentTok{\# Combine top biomarkers into a penalized panel}
\NormalTok{panel }\OtherTok{\textless{}{-}} \FunctionTok{fit\_auc\_panel}\NormalTok{(X[, }\FunctionTok{head}\NormalTok{(cv}\SpecialCharTok{$}\NormalTok{results}\SpecialCharTok{$}\NormalTok{biomarker, }\DecValTok{4}\NormalTok{)], y,}
  \AttributeTok{method =} \StringTok{"ridge"}\NormalTok{, }\AttributeTok{n\_folds =} \DecValTok{5}\NormalTok{)}

\CommentTok{\# Sample size for 80\% power to detect AUC = 0.70}
\FunctionTok{power\_auc\_matrix}\NormalTok{(}
  \AttributeTok{target\_aucs =} \FloatTok{0.70}\NormalTok{, }\AttributeTok{power =} \FloatTok{0.80}\NormalTok{,}
  \AttributeTok{prevalence =} \FloatTok{0.3}\NormalTok{, }\AttributeTok{alpha =} \FloatTok{0.05}
\NormalTok{)}
\end{Highlighting}
\end{Shaded}

\section{Benchmarking}\label{benchmarking}

\textbf{Screening throughput.} On a desktop workstation (AMD Ryzen 3.5
GHz, 32 GB RAM, R 4.5.1, Windows 11), aucmat screens 100, 1,000, and
10,000 biomarkers (n = 500, DeLong inference) in 0.54, 5.02, and 59.76
seconds respectively. The equivalent per-column pROC loop
(\texttt{pROC::roc()} then \texttt{pROC::ci.auc()} with
\texttt{method\ =\ "delong"}) on the same data requires 0.68, 6.42, and
an estimated 64 seconds -- a speedup of approximately 1.1--1.3× at
moderate \(p\), increasing at larger \(p\) as aucmat's fixed per-call
overhead is amortized over more columns. The speed advantage comes from
avoiding repeated outcome encoding and from column-wise rank
computation. At small \(p\) (\textless{} 100), the difference is
negligible; for wide matrices (\(p > 1{,}000\)), the per-column overhead
of \texttt{pROC::roc()} and \texttt{pROC::ci.auc()} becomes meaningful.

\textbf{Simulation calibration.} Across 27 simulation conditions (n =
100, 300, 500; prevalence averaged over 0.2, 0.3, 0.5; target AUC =
0.65, 0.75, 0.85; 200 replicates each; averaged over exchangeable and
AR(1) correlation structures with \(\rho = 0\) and \(\rho = 0.3\)),
DeLong 95\% confidence interval coverage ranged from 0.898 (at n = 100,
AUC = 0.65) to 0.957 (at n = 500, AUC = 0.65), converging to the nominal
0.95 as sample size increased. Power exceeded 0.85 for AUC \(\geq\) 0.65
at n = 100 and reached 1.0 for AUC \(\geq\) 0.75 at n \(\geq\) 300. The
coverage is slightly below nominal for small \(n\) and extreme AUC
values, consistent with the known finite-sample behavior of the DeLong
variance estimator.

\begin{longtable}[]{@{}lrrrr@{}}
\caption{Empirical coverage of DeLong 95\% confidence intervals,
averaged across prevalence (0.2, 0.3, 0.5) and correlation structures
(exchangeable, AR1) with 200 replicates each.}\tabularnewline
\toprule\noalign{}
& AUC & n=100 & n=300 & n=500 \\
\midrule\noalign{}
\endfirsthead
\toprule\noalign{}
& AUC & n=100 & n=300 & n=500 \\
\midrule\noalign{}
\endhead
\bottomrule\noalign{}
\endlastfoot
1 & 0.65 & 0.898 & 0.937 & 0.957 \\
4 & 0.75 & 0.915 & 0.943 & 0.933 \\
7 & 0.85 & 0.883 & 0.923 & 0.905 \\
\end{longtable}

\section{Comparison with existing R
packages}\label{comparison-with-existing-r-packages}

Table \ref{features} compares aucmat (v0.2.0) with five established R
packages for ROC analysis and biomarker evaluation. Features are
verified against each package's CRAN documentation and published
literature as of July 2026.

\begin{longtable}[]{@{}
  >{\raggedright\arraybackslash}p{(\linewidth - 12\tabcolsep) * \real{0.5125}}
  >{\raggedright\arraybackslash}p{(\linewidth - 12\tabcolsep) * \real{0.0875}}
  >{\raggedright\arraybackslash}p{(\linewidth - 12\tabcolsep) * \real{0.0625}}
  >{\raggedright\arraybackslash}p{(\linewidth - 12\tabcolsep) * \real{0.1000}}
  >{\raggedright\arraybackslash}p{(\linewidth - 12\tabcolsep) * \real{0.0625}}
  >{\raggedright\arraybackslash}p{(\linewidth - 12\tabcolsep) * \real{0.0875}}
  >{\raggedright\arraybackslash}p{(\linewidth - 12\tabcolsep) * \real{0.0875}}@{}}
\caption{Feature comparison of R packages for ROC analysis and biomarker
screening. \textbackslash texttt\{+\} = supported;
\textbackslash texttt\{-\} = not supported or requires manual
implementation. Verified against CRAN documentation as of July 2026.
pROC supports unpaired comparisons via
\textbackslash texttt\{roc.test(paired=FALSE)\}; aucmat requires common
subjects for all comparisons.}\tabularnewline
\toprule\noalign{}
\begin{minipage}[b]{\linewidth}\raggedright
Feature
\end{minipage} & \begin{minipage}[b]{\linewidth}\raggedright
aucmat
\end{minipage} & \begin{minipage}[b]{\linewidth}\raggedright
pROC
\end{minipage} & \begin{minipage}[b]{\linewidth}\raggedright
precrec
\end{minipage} & \begin{minipage}[b]{\linewidth}\raggedright
ROCR
\end{minipage} & \begin{minipage}[b]{\linewidth}\raggedright
colAUC
\end{minipage} & \begin{minipage}[b]{\linewidth}\raggedright
dtComb
\end{minipage} \\
\midrule\noalign{}
\endfirsthead
\toprule\noalign{}
\begin{minipage}[b]{\linewidth}\raggedright
Feature
\end{minipage} & \begin{minipage}[b]{\linewidth}\raggedright
aucmat
\end{minipage} & \begin{minipage}[b]{\linewidth}\raggedright
pROC
\end{minipage} & \begin{minipage}[b]{\linewidth}\raggedright
precrec
\end{minipage} & \begin{minipage}[b]{\linewidth}\raggedright
ROCR
\end{minipage} & \begin{minipage}[b]{\linewidth}\raggedright
colAUC
\end{minipage} & \begin{minipage}[b]{\linewidth}\raggedright
dtComb
\end{minipage} \\
\midrule\noalign{}
\endhead
\bottomrule\noalign{}
\endlastfoot
Matrix-first screening (single call) & + & - & - & - & + & - \\
Direction-preserving AUC (\textless0.5 reported) & + & - & - & - & - &
- \\
DeLong confidence intervals & + & + & - & - & - & - \\
Stratified bootstrap CI & + & + & - & - & - & + \\
Multiplicity adjustment (BH/Holm/Bonf) & + & - & - & - & - & - \\
One-sided alternative tests & + & + & - & - & - & - \\
Paired biomarker comparisons & + & + & - & - & - & + \\
Non-inferiority / equivalence (TOST) & + & - & - & - & - & - \\
Omnibus Wald test (3+ correlated AUCs) & + & - & - & - & - & - \\
Bootstrap rank-stability analysis & + & - & - & - & - & - \\
Partial AUC & + & + & + & - & - & - \\
Multivariate simulation engine & + & - & - & - & - & - \\
Simulation calibration diagnostics & + & - & - & - & - & - \\
Cross-validated AUC screening & + & - & - & - & - & + \\
Penalized panel models (ridge/lasso) & + & - & - & - & - & + \\
Hurdle-AUC (zero-inflated data) & + & - & - & - & - & - \\
Precision-Recall curves & + & - & + & - & - & - \\
Power / sample size computation & + & + & - & - & - & - \\
Multiclass AUC (Hand-Till) & + & + & - & - & + & - \\
ROC smoothing (binormal/density) & + & + & - & - & - & - \\
Group-stratified screening & + & - & - & - & - & - \\
\end{longtable}

\textbf{What aucmat uniquely provides.} Five capabilities are, to our
knowledge, not available in any other single R package: (i) matrix-first
screening with direction-preserving AUC and multiplicity adjustment in
one call; (ii) non-inferiority and equivalence (TOST) tests for paired
AUC comparisons; (iii) an omnibus Wald test for three or more correlated
AUCs via the joint DeLong covariance matrix; (iv) bootstrap
rank-stability analysis with top-\(k\) selection probabilities and
pairwise co-selection frequencies; and (v) a two-stage Hurdle-AUC model
for zero-inflated data that recovers discriminative signal where
standard AUC fails.

\textbf{When to use what.} For a single biomarker with smoothed ROC
curves, pROC remains the gold standard, offering binormal and density
smoothing, partial AUC with bootstrap inference, and Venkatraman's
permutation test for whole-curve comparison \citep{Venkatraman2000}. For
speed on extremely wide matrices (\textgreater100k columns), colAUC's
Wilcoxon-based column-wise computation is recommended, though it
provides no inference. For precision-recall analysis of highly
imbalanced data, precrec's C++ backend offers fast and accurate PR
curves with confidence bands for multiple test sets. For combining two
biomarkers into an optimal composite score, dtComb's library of 142
methods covers virtually all published linear, non-linear, mathematical,
and machine-learning combination approaches. For end-to-end screening of
hundreds to tens of thousands of biomarkers with full statistical
inference, multiplicity adjustment, comparison testing, stability
assessment, validation, and study planning in a single coherent
workflow, aucmat is the recommended tool.

\section{Discussion}\label{discussion}

aucmat provides a unified, statistically principled framework for
biomarker screening against binary outcomes in R. Its matrix-first
design reduces the code burden for high-dimensional ROC analysis from
dozens of bespoke scripting steps to a handful of function calls with
consistent S3 interfaces. Cross-validated screening and penalized panel
models address same-data optimism, while power analysis tools support
prospective study design.

\textbf{Limitations.} The current version (0.2.0) is restricted to
binary outcomes for the primary screening workflow (multiclass is
available experimentally via \texttt{aucmat\_multiclass()}). The DeLong
variance estimator requires common subjects for paired comparisons;
unpaired comparisons between independent cohorts are not yet implemented
(pROC provides this via \texttt{roc.test(paired=FALSE)}). Partial AUC
computation is implemented as an internal engine
(\texttt{R/partial-auc.R}) but not yet exposed through the primary
\texttt{aucmat()} interface. The package targets in-memory dense
matrices; sparse and on-disk backends are deferred. The DeLong variance
estimator shows slightly liberal coverage at small sample sizes
(coverage \(\approx\) 0.90 at n = 100 for AUC = 0.65), consistent with
its known finite-sample behavior; bootstrap intervals are recommended
when n \textless{} 100 per class.

\textbf{Future work.} The development roadmap targets four major
extensions: version 0.3 will add multivariable penalized panel models
with honest nested cross-validation to prevent information leakage
across preprocessing, feature selection, and tuning steps; version 0.4
will extend the framework to multiclass outcomes with one-versus-rest,
one-versus-one, and Hand-Till estimators; version 0.5 will incorporate
cumulative/dynamic time-dependent AUCs for censored survival outcomes
\citep{Heagerty2000, Blanche2013}; and version 0.6 will introduce
sparse-matrix and on-disk backends for genomics-scale data exceeding
available RAM.

\section{Acknowledgements}\label{acknowledgements}

We thank Xavier Robin and colleagues for pROC, which provided both the
statistical foundation and the ROC curve objects that aucmat builds
upon. Hadley Wickham's ggplot2 and the mvtnorm authors (Genz, Bretz, et
al.) are gratefully acknowledged for the visualization and simulation
backends. The Sun \& Xu (2014) fast DeLong algorithm implemented in pROC
since version 1.9 informs aucmat's DeLong variance computation.

Thank KUHNE foundation for financial support.

\renewcommand\refname{References}
\bibliography{refs.bib}

\end{document}
